\documentclass[11pt,a4paper]{article}

\usepackage[T1]{fontenc}
\usepackage{lmodern}
\usepackage{microtype}
\usepackage[a4paper,margin=27mm]{geometry}
\usepackage{amsmath,amssymb,amsthm,mathtools}
\usepackage{booktabs,array}
\usepackage{enumitem}
\usepackage{xcolor}
\usepackage[colorlinks=true,linkcolor=blue!55!black,citecolor=blue!55!black,urlcolor=blue!55!black]{hyperref}
\hypersetup{pdftitle={A Local Matrix-Integrating-Factor Equivalence for Kerr--Newman Perturbation Systems},
pdfauthor={He Wang; GPT-6.1 Sol; Claude Code Opus 5.5},pdfsubject={Local equivalence of coupled first-order perturbation systems}}

\allowdisplaybreaks[1]
\setlist{leftmargin=*,itemsep=2pt,topsep=3pt}
\numberwithin{equation}{section}

\newtheorem{theorem}{Theorem}[section]
\newtheorem{proposition}[theorem]{Proposition}
\newtheorem{lemma}[theorem]{Lemma}
\newtheorem{corollary}[theorem]{Corollary}
\theoremstyle{definition}
\newtheorem{definition}[theorem]{Definition}
\theoremstyle{remark}
\newtheorem{remark}[theorem]{Remark}

\newcommand{\R}{\mathbb R}
\newcommand{\C}{\mathbb C}
\newcommand{\I}{\mathrm i}
\newcommand{\tr}{\operatorname{tr}}
\newcommand{\adj}{\operatorname{adj}}
\newcommand{\Uext}{\mathcal U}
\newcommand{\D}{\mathcal D}
\newcommand{\Lop}{\mathcal L}
\newcommand{\Res}{\mathcal R}
\newcommand{\mat}[1]{\begin{pmatrix}#1\end{pmatrix}}
\newcommand{\cE}{\mathcal E}

\title{\textbf{A Local Matrix-Integrating-Factor Equivalence\\
for Kerr--Newman Perturbation Systems}}
\author{He Wang\thanks{ORCID: \href{https://orcid.org/0000-0002-1353-391X}{0000-0002-1353-391X}.}
\and GPT-6.1 Sol\thanks{AI contributor (OpenAI): assistance with mathematical exposition, drafting, and manuscript checks.}
\and Claude Code Opus 5.5\thanks{AI contributor (Anthropic).}}
\date{Draft, 29 October 2026}

\begin{document}
\maketitle

\begin{abstract}
We establish a local equivalence between two coupled first-order systems
for perturbations of a subextremal Kerr--Newman black hole.  On the exterior
coordinate patch \(r>r_+\), \(0<\theta<\pi\), we construct chiral
two-component fields from each of the positive- and negative-spin
four-field systems.  A parameter-dependent matrix \(X(r,\theta)\) has
constant nonzero determinant and satisfies a first-order integrating-factor
equation.  These identities remove the matrix connection and yield an
explicit transformed system.  We prove both directions of the equivalence,
including the inverse field maps and all nonvanishing conditions needed in
the calculation.  The result holds pointwise for every real frequency and
azimuthal parameter and arbitrary smooth complex fields on the patch.  It
does not assert a separation theorem, a global boundary-value result, or
mode stability; those are distinct analytic questions.
\end{abstract}

\noindent\textbf{Keywords:} Kerr--Newman; Einstein--Maxwell perturbations;
matrix integrating factor; chiral variables; first-order systems.

\section{Introduction and main result}

The perturbation equations of a charged rotating black hole couple
spin-two and spin-one radiative degrees of freedom.  Hintz's
first-order formulation \cite{Hintz2026} treats modes with dependence
\(\exp(-\I\omega t+\I m\phi)\) on the stationary and azimuthal coordinates.
In the formulation considered here, each spin orientation is described by four complex
fields.  A particular chiral recombination packages those fields into two
complex two-vectors.  The resulting system contains a matrix connection
\(A\).  The key identity is that a second matrix \(X\) differentiates
according to that same connection.  Multiplying the chiral variables by
\(\adj(X)\) and \(X^\dagger\) then removes the connection from the
derivatives.  This paper gives a detailed verification of that local
transformation, with explicit residual identities for both spin signs.
The equations and transformation originate in \cite{Hintz2026}; our
purpose is to expose the complete chain of identities and the hypotheses
under which both directions are valid.

We retain the names system (1), system (4), and system (7)
from \cite{Hintz2026} to distinguish the original, chiral, and
transformed formulations.  Displayed equations in this paper have their
own section-based numbering.  We write \(\rho=r+\I a\cos\theta\) for
the factor denoted \(\bar\rho\) in \cite{Hintz2026}, and use
\(X^\dagger\) for complex conjugate transpose; the source writes this
matrix as \(X^*\).  The spin sign \(\varsigma=\pm1\) labels the two
four-field sectors, each of which contains the coupled gravitational and
electromagnetic components.

We state precisely the local claim proved below.  Let
\begin{equation}
 M>0,\qquad Q\neq0,\qquad a^2+Q^2<M^2,
\label{eq:parameters}
\end{equation}
and define
\begin{equation}
 \Delta(r)=r^2-2Mr+a^2+Q^2,\qquad
 r_+=M+\sqrt{M^2-a^2-Q^2}.
\label{eq:delta}
\end{equation}
In geometrized units, \(M\) is the mass, \(Ma\) the angular momentum,
and \(Q\) the electric charge.
The domain is the open exterior patch
\begin{equation}
 \Uext=\{(r,\theta)\in\R^2:r>r_+,\ 0<\theta<\pi\}.
\label{eq:domain}
\end{equation}
For real \(\omega,m\), let \(F=(\psi_2,\psi_1,k,s)\) denote a
positive-spin field and \(G=(\psi_{-2},\psi_{-1},n,\ell)\) a
negative-spin field.

\begin{theorem}[Local equivalence of systems (1), (4), and (7)]
\label{thm:main}
Suppose the parameters satisfy \eqref{eq:parameters}.  For every
\(\omega,m\in\R\), every \(C^\infty\) positive-spin field \(F\) on
\(\Uext\), and each \(x\in\Uext\), the four positive-spin equations of
system (1) hold at \(x\) if and only if the transformed fields
\[
 z_+=\adj(X)y_+,\qquad z_-=X^\dagger y_-
\]
satisfy system (7) at \(x\).  Here \(y_\pm\) are the chiral combinations
defined in \eqref{eq:chiral-positive}, and \(X\) is the matrix in
\eqref{eq:X}.

The same equivalence holds for every \(C^\infty\) negative-spin field
\(G\), using the negative-spin chiral combinations
\eqref{eq:chiral-negative}.  In both sectors the result is pointwise on
all of \(\Uext\), and hence also holds when the pointwise equivalence is
quantified over the entire patch.
\end{theorem}

The statement concerns equivalence of differential equations under an
explicit field transformation.  It does not claim existence or uniqueness
of solutions with prescribed boundary data.  In particular, it is not a
global mode-stability theorem.

\section{Geometry, operators, and the original system}
\label{sec:original}

\subsection{Exterior identities}

On \(\Uext\), the assumptions imply
\begin{equation}
 \Delta>0,\qquad \sin\theta>0,\qquad Q^2-M^2\neq0.
\label{eq:positivity}
\end{equation}
Indeed, writing \(r_-=M-\sqrt{M^2-a^2-Q^2}\), one has
\(\Delta=(r-r_+)(r-r_-)\), so \(r>r_+\) implies \(\Delta>0\).
Write
\begin{equation}
 \rho=r+\I a\cos\theta,\qquad
 \rho^*=r-\I a\cos\theta=\overline{\rho}.
\label{eq:rho}
\end{equation}
Since \(r>r_+>0\), neither \(\rho\) nor \(\rho^*\) vanishes.  Thus all
expressions below with denominators \(\Delta,\sin\theta,\rho,\rho^*\), or
\(Q\) are defined throughout the stated domain.  These are the only
domain-level exclusions used in the algebraic transformations.

\subsection{Physical meaning of the four-field variables}
\label{sec:physical-fields}

For orientation, we record the field conventions of
\cite{Hintz2026}.  In its Newman--Penrose tetrad and gauge, the
linearized Weyl scalars are rescaled as
\begin{align}
 \psi_2&=\Psi_0^{(1)},&
 \psi_1&=\sqrt2\,\rho^*\Psi_1^{(1)},&
 \psi_{-1}&=\frac{(\rho^*)^3}{\sqrt2}\Psi_3^{(1)},&
 \psi_{-2}&=(\rho^*)^4\Psi_4^{(1)}.
 \label{eq:physical-weyl}
\end{align}
The remaining fields are rescaled linearized spin coefficients:
\begin{equation}
 k=\frac{\kappa^{(1)}}{\sqrt2\,(\rho^*)^2},\qquad
 s=\frac{\rho\,\sigma^{(1)}}{(\rho^*)^2},\qquad
 n=\frac{\rho\rho^*}{\sqrt2}\nu^{(1)},\qquad
 \ell=\frac{\rho^*}{2}\lambda^{(1)}.
 \label{eq:physical-spin}
\end{equation}
The superscript \((1)\) here denotes linearization.  The lower-case
subscripts on \(\psi_{\pm2},\psi_{\pm1}\) label spin weights; upper-case
subscripts follow the Newman--Penrose numbering of the Weyl scalars.
Each quartet combines curvature and connection data, with the
Einstein--Maxwell coupling encoded in the first-order system below.
The dictionary \eqref{eq:physical-weyl}--\eqref{eq:physical-spin}
provides physical context from \cite{Hintz2026}; our proof starts with
that system and treats its four fields as arbitrary smooth functions.

\subsection{Mode operators}

For real mode parameters \(\omega,m\), set
\begin{equation}
 K(r)=-(r^2+a^2)\omega+am,\qquad
 \widehat Q(\theta)=-a\omega\sin\theta+\frac{m}{\sin\theta}.
\label{eq:KQ}
\end{equation}
For a complex-valued field \(f=f(r,\theta)\), define
\begin{align}
 \D_j f&=\partial_r f+
 \left(\frac{\I K}{\Delta}+j\frac{\Delta'}{\Delta}\right)f,
&
 \D_j^\dagger f&=\partial_r f+
 \left(-\frac{\I K}{\Delta}+j\frac{\Delta'}{\Delta}\right)f,
\label{eq:Doperators}\\
 \Lop_j f&=\partial_\theta f+
 \left(\widehat Q+j\cot\theta\right)f,
&
 \Lop_j^\dagger f&=\partial_\theta f+
 \left(-\widehat Q+j\cot\theta\right)f.
\label{eq:Loperators}
\end{align}
The daggers on \(\D_j^\dagger\) and \(\Lop_j^\dagger\) specify the
displayed sign changes; no Hilbert-space adjoint or measure is being
introduced.  In contrast, a dagger on a matrix always means its complex
conjugate transpose.
The complex first-order derivatives used in the proof are
\begin{equation}
 D_-=\sqrt{\Delta}\,\partial_r-\I\,\partial_\theta,\qquad
 D_+=\sqrt{\Delta}\,\partial_r+\I\,\partial_\theta.
\label{eq:Dpm}
\end{equation}
They act componentwise on vector- and matrix-valued fields.

\subsection{Positive-spin residuals}

Set
\begin{equation}
 F_\pm=3\left(M-\frac{Q^2}{\rho}\right)
 \mathbin{\pm}\frac{Q^2\rho^*}{\rho^2}.
\label{eq:Fpm}
\end{equation}
The positive-spin equations are the vanishing of the following four
residuals:
\begin{align}
 \Res_{1a}={}&
 \left(\Lop_2-\frac{3\I a\sin\theta}{\rho^*}\right)\psi_2
 -\left(\D_0+\frac{3}{\rho^*}\right)\psi_1+2kF_+,
\label{eq:R1a}\\
 \Res_{1b}={}&
 \Delta\left(\D_2^\dagger-\frac{3}{\rho^*}\right)\psi_2
 +\left(\Lop_{-1}^\dagger+\frac{3\I a\sin\theta}{\rho^*}\right)\psi_1
 -2sF_-,
\label{eq:R1b}\\
 \Res_{1c}={}&
 \left(\D_0+\frac{3}{\rho^*}\right)s
 -\left(\Lop_{-1}^\dagger+\frac{3\I a\sin\theta}{\rho^*}\right)k
 -\frac{\rho}{(\rho^*)^2}\psi_2,
\label{eq:R1c}\\
 \Res_{1d}={}&
 \Delta\left(\D_2^\dagger-\frac{3}{\rho^*}\right)k
 +\left(\Lop_2-\frac{3\I a\sin\theta}{\rho^*}\right)s
 -\frac{2\rho}{(\rho^*)^2}\psi_1.
\label{eq:R1d}
\end{align}
Thus system (1) is \(\Res_{1a}=\Res_{1b}=\Res_{1c}=\Res_{1d}=0\).

For the negative-spin system, apply the replacements
\begin{equation}
 \psi_2\mapsto\psi_{-2},\qquad
 \psi_1\mapsto-\psi_{-1},\qquad
 k\mapsto-n,\qquad s\mapsto\ell,
\label{eq:negativefields}
\end{equation}
and
\begin{equation}
 \Lop_{-1}^\dagger\mapsto\Lop_{-1},\quad
 \Lop_2\mapsto\Lop_2^\dagger,\quad
 \D_0+\frac3{\rho^*}\mapsto
 \Delta\left(\D_{-1}^\dagger+\frac3{\rho^*}\right),\quad
 \Delta\left(\D_2^\dagger-\frac3{\rho^*}\right)
 \mapsto\D_0-\frac3{\rho^*}.
\label{eq:negativeoperators}
\end{equation}
Denote the resulting four residuals by
\(\Res^-_{1a},\ldots,\Res^-_{1d}\).  System (1)\({}^-\) is the
simultaneous vanishing of these four functions.
\begin{align}
 \Res^-_{1a}={}&
 \left(\Lop_2^\dagger-\frac{3\I a\sin\theta}{\rho^*}\right)\psi_{-2}
 -\Delta\left(\D_{-1}^\dagger+\frac{3}{\rho^*}\right)(-\psi_{-1})
 -2nF_+,\label{eq:Rminus1a}\\
 \Res^-_{1b}={}&
 \left(\D_0-\frac{3}{\rho^*}\right)\psi_{-2}
 +\left(\Lop_{-1}+\frac{3\I a\sin\theta}{\rho^*}\right)(-\psi_{-1})
 -2\ell F_-,\label{eq:Rminus1b}\\
 \Res^-_{1c}={}&
 \Delta\left(\D_{-1}^\dagger+\frac{3}{\rho^*}\right)\ell
 -\left(\Lop_{-1}+\frac{3\I a\sin\theta}{\rho^*}\right)(-n)
 -\frac{\rho}{(\rho^*)^2}\psi_{-2},\label{eq:Rminus1c}\\
 \Res^-_{1d}={}&
 \left(\D_0-\frac{3}{\rho^*}\right)(-n)
 +\left(\Lop_2^\dagger-\frac{3\I a\sin\theta}{\rho^*}\right)\ell
 -\frac{2\rho}{(\rho^*)^2}(-\psi_{-1}).\label{eq:Rminus1d}
\end{align}

\section{Chiral recombination and system (4)}
\label{sec:chiral}

\subsection{Chiral fields}

Define the positive prefactors
\begin{equation}
 P_+=\Delta^{5/4}\sqrt{\sin\theta},\qquad
 P_-=\Delta^{-3/4}\sqrt{\sin\theta}.
\label{eq:prefactors}
\end{equation}
For \(\varepsilon=\pm1\), the positive-spin chiral vectors are
\begin{equation}
 y_\varepsilon^{(+)}
 =P_+\mat{
 \psi_2+\varepsilon\I\Delta^{-1/2}\psi_1\\
 -2Q(k+\varepsilon\I\Delta^{-1/2}s)}.
\label{eq:chiral-positive}
\end{equation}
The negative-spin chiral vectors are
\begin{equation}
 y_\varepsilon^{(-)}
 =P_-\mat{
 \psi_{-2}-\varepsilon\I\Delta^{1/2}\psi_{-1}\\
 2Q(n-\varepsilon\I\Delta^{1/2}\ell)}.
\label{eq:chiral-negative}
\end{equation}
The superscript labels the spin sector and the subscript labels the two
chiral combinations.  Positivity of \(P_\pm\) on \(\Uext\) follows from
\eqref{eq:positivity}.

\subsection{The coefficient matrices}

For a spin sign \(\varsigma\in\{+1,-1\}\), define
\begin{align}
 c^{(\varsigma)}
 &=\left(\frac{\varsigma\I K}{\Delta}
 +\frac3{\rho^*}-\frac{3\Delta'}{4\Delta}\right)\sqrt{\Delta},
\label{eq:c}\\
 b^{(\varsigma)}
 &=-\varsigma\left(-a\omega\sin\theta+\frac{m}{\sin\theta}\right)
 -\frac32\frac{\cos\theta}{\sin\theta}
 +\frac{3\I a\sin\theta}{\rho^*},\nonumber\\
 B^{(\varsigma)}
 &=\I\mat{
 b^{(\varsigma)}&
 \displaystyle\frac{3(M\rho-Q^2)}{Q\rho}\\[4pt]
 \displaystyle\frac{3Q\rho}{(\rho^*)^2}&-b^{(\varsigma)}
},
\label{eq:B}\\
 A&=\mat{
 0&\displaystyle-\I Q\frac{\rho^*}{\rho^2}\\[4pt]
 \displaystyle\I Q\frac{\rho}{(\rho^*)^2}&0}.
\label{eq:A}
\end{align}
In particular,
\begin{equation}
 \tr B^{(\varsigma)}=0,\qquad A^\dagger=A.
\label{eq:trace-hermitian}
\end{equation}

\subsection{Intermediate chiral equations}

System (4) is the pair
\begin{align}
 D_-y_+^{(\varsigma)}-Ay_+^{(\varsigma)}
 &=\left(c^{(\varsigma)}I-B^{(\varsigma)}\right)y_-^{(\varsigma)},
\label{eq:S4minus}\\
 D_+y_-^{(\varsigma)}+Ay_-^{(\varsigma)}
 &=\left(c^{(\varsigma)}I+B^{(\varsigma)}\right)y_+^{(\varsigma)}.
\label{eq:S4plus}
\end{align}
The \(I\) multiplying \(c^{(\varsigma)}\) is the \(2\times2\) identity.

\begin{proposition}[Exact residual identities for chiralization]
\label{prop:residual}
Let \(\cE_-^{(\pm)}\) be the left side of \eqref{eq:S4minus} and
\(\cE_+^{(\pm)}\) the left side of \eqref{eq:S4plus} after moving the
right side to the left.  For positive spin, their components are
\begin{align}
 \cE_{-,0}^{(+)}
 &=P_+\left(-\I\Res_{1a}
 +\frac{\Res_{1b}}{\sqrt\Delta}\right),&
 \cE_{-,1}^{(+)}
 &=P_+\left[-2Q\left(\I\Res_{1c}
 +\frac{\Res_{1d}}{\sqrt\Delta}\right)\right],
\label{eq:positive-residuals}\\
 \cE_{+,0}^{(+)}
 &=P_+\left(\I\Res_{1a}
 +\frac{\Res_{1b}}{\sqrt\Delta}\right),&
 \cE_{+,1}^{(+)}
 &=P_+\left[-2Q\left(-\I\Res_{1c}
 +\frac{\Res_{1d}}{\sqrt\Delta}\right)\right].
\label{eq:positive-residuals-plus}
\end{align}
For negative spin, the corresponding identities are
\begin{align}
 \cE_{-,0}^{(-)}
 &=P_-\left(-\I\Res^-_{1a}
 +\sqrt\Delta\,\Res^-_{1b}\right),&
 \cE_{-,1}^{(-)}
 &=P_-\left[2Q\left(-\I\Res^-_{1c}
 -\sqrt\Delta\,\Res^-_{1d}\right)\right],
\label{eq:negative-residuals}\\
 \cE_{+,0}^{(-)}
 &=P_-\left(\I\Res^-_{1a}
 +\sqrt\Delta\,\Res^-_{1b}\right),&
 \cE_{+,1}^{(-)}
 &=P_-\left[2Q\left(\I\Res^-_{1c}
 -\sqrt\Delta\,\Res^-_{1d}\right)\right].
\label{eq:negative-residuals-plus}
\end{align}
Consequently, in either spin sector system (1) is equivalent at each
\(x\in\Uext\) to system (4).
\end{proposition}

\begin{proof}
The derivative of a prefactor is computed from
\(\Delta'=2(r-M)\), \(\partial_r\sqrt\Delta=(r-M)/\sqrt\Delta\),
and \(\partial_\theta\sqrt{\sin\theta}
=\cos\theta/(2\sqrt{\sin\theta})\).  It follows that
\[
 \frac{D_-P_+}{P_+}
 =\frac{5\Delta'}{4\sqrt\Delta}-\frac{\I}{2}\cot\theta,\qquad
 \frac{D_+P_+}{P_+}
 =\frac{5\Delta'}{4\sqrt\Delta}+\frac{\I}{2}\cot\theta,
\]
and the same identities hold for \(P_-\) with \(5/4\) replaced by
\(-3/4\).  The remaining derivatives needed in the four components are
\[
 \partial_r\Delta^{-1/2}
 =-\frac{\Delta'}{2\Delta^{3/2}},\qquad
 \partial_r\Delta^{1/2}
 =\frac{\Delta'}{2\sqrt\Delta}.
\]
Apply the product rule to each component of
\eqref{eq:chiral-positive} and \eqref{eq:chiral-negative}.  Substitute
the definitions of \(D_\pm\), \(A\), \(B^{(\varsigma)}\), and
\(c^{(\varsigma)}\).  The coefficients of each field derivative match
the corresponding terms in \eqref{eq:R1a}--\eqref{eq:R1d} and their
negative-spin replacements.  The zero-order terms agree by the
definitions of \(b,c,F_\pm\) and the off-diagonal entries of \(A,B\).
Collecting terms gives the eight displayed component identities.
Appendix~\ref{app:component} gives both a worked component and a uniform
operator calculation covering all eight components.

If all original residuals vanish, the displayed formulas give system (4).
For the converse, divide each component by its nonzero prefactor.
With \(S=\sqrt\Delta\), the coefficient matrices for the positive-spin
residual pairs \((\Res_{1a},\Res_{1b})\) and
\((\Res_{1c},\Res_{1d})\) are, respectively,
\[
 \mat{-\I&S^{-1}\\ \I&S^{-1}},
 \qquad \mat{\I&S^{-1}\\ -\I&S^{-1}}.
\]
Their determinants are \(-2\I/S\) and \(2\I/S\).  The corresponding
negative-spin matrices are
\[
 \mat{-\I&S\\ \I&S},
 \qquad \mat{-\I&-S\\ \I&-S},
\]
with determinants \(-2\I S\) and \(2\I S\).  All four determinants
are nonzero.  Therefore vanishing of the four chiral residual components
forces vanishing of the four original residuals in either sector.
\end{proof}

\begin{remark}[Zeros of the original coupling]
\label{rem:coupling-zeros}
The proof divides by \(Q,P_\pm,\sqrt\Delta\), and fixed nonzero
numerical factors, but never by \(F_+\) or \(F_-\).  Consequently
their zero sets cause no exception to the equivalence.  Such zeros
can occur in the physical exterior: when \(a=0\),
\(F_-=3M-4Q^2/r\).  For example, \(M=1\), \(Q^2=63/64\) gives
\(r_+=9/8\) and \(F_-(21/16)=0\), with \(21/16>9/8\).
\end{remark}

\section{The matrix integrating factor}
\label{sec:integrating}

\subsection{Definition and constant determinant}

Define
\begin{equation}
 X(r,\theta)=
 \mat{
 \displaystyle Q\frac{\sqrt\Delta+a\sin\theta}{\rho}
 &\displaystyle-\I\left(M-\frac{Q^2}{\rho}\right)\\[7pt]
 \displaystyle\I\left(M-\frac{Q^2}{\rho^*}\right)
 &\displaystyle Q\frac{\sqrt\Delta-a\sin\theta}{\rho^*}
}.
\label{eq:X}
\end{equation}
Each entry is smooth on \(\Uext\).  To compute the determinant, expand the
two products in \(X_{00}X_{11}-X_{01}X_{10}\) and multiply by
\(\rho\rho^*\), which is nonzero on the domain.  Use
\[
 \rho\rho^*=r^2+a^2\cos^2\theta,\quad
 (\sqrt\Delta)^2=r^2-2Mr+a^2+Q^2,\quad
 \sin^2\theta+\cos^2\theta=1.
\]
More explicitly, the numerator is
\begin{align*}
 \rho\rho^*\det X
 &=Q^2(\Delta-a^2\sin^2\theta)
   -(M\rho-Q^2)(M\rho^*-Q^2)\\
 &=Q^2(r^2-2Mr+a^2\cos^2\theta+Q^2)\\
 &\quad-M^2(r^2+a^2\cos^2\theta)+2MrQ^2-Q^4\\
 &=(Q^2-M^2)(r^2+a^2\cos^2\theta).
\end{align*}
Canceling the nonzero factor \(\rho\rho^*\), we obtain
\begin{equation}
 \det X=Q^2-M^2.
\label{eq:detX}
\end{equation}
By \eqref{eq:parameters}, this constant is nonzero, so \(X\) is
invertible everywhere on \(\Uext\).  In dimension two,
\begin{equation}
 \adj(X)=(\det X)X^{-1}.
\label{eq:adjX}
\end{equation}

\subsection{Differentiating \texorpdfstring{\(X\)}{X}}

Put \(S=\sqrt\Delta\), \(s=\sin\theta\), and \(c_\theta=\cos\theta\).
The four entries of \(A\) and \(X\) are explicit rational functions of
\(\rho,\rho^*\), \(S\), and \(s\).  Their derivatives are controlled by
\begin{align}
 \Delta'&=2(r-M),&
 \partial_r\sqrt\Delta&=\frac{r-M}{\sqrt\Delta},\label{eq:delta-derivs}\\
 \partial_r\rho&=\partial_r\rho^*=1,&
 \partial_\theta\rho&=-\I a\sin\theta,&
 \partial_\theta\rho^*&=\I a\sin\theta.\label{eq:rho-derivs}
\end{align}
In particular,
\[
 D_-\rho=S-as,\qquad D_-\rho^*=S+as,\qquad
 D_-(S\pm as)=r-M\mp\I a c_\theta.
\]
The quotient rule gives the diagonal derivatives explicitly:
\begin{align}
 D_-X_{00}
 &=Q\frac{(r-M-\I a c_\theta)\rho-(S+as)(S-as)}{\rho^2}\nonumber\\*
 &=Q\frac{(\rho^*-M)\rho-(\Delta-a^2s^2)}{\rho^2}
 =Q\frac{M\rho^*-Q^2}{\rho^2},\label{eq:X00deriv}\\
 D_-X_{11}
 &=Q\frac{(r-M+\I a c_\theta)\rho^*-(S-as)(S+as)}
 {(\rho^*)^2}\nonumber\\*
 &=Q\frac{(\rho-M)\rho^*-(\Delta-a^2s^2)}{(\rho^*)^2}
 =Q\frac{M\rho-Q^2}{(\rho^*)^2}.\label{eq:X11deriv}
\end{align}
In the last equalities use
\(\rho\rho^*=r^2+a^2c_\theta^2\),
\(s^2+c_\theta^2=1\), and the defining polynomial for \(\Delta\).
For the off-diagonal entries write
\(X_{01}=-\I M+\I Q^2/\rho\) and
\(X_{10}=\I M-\I Q^2/\rho^*\); differentiating the reciprocals gives
\begin{equation}
 D_-X_{01}=-\I Q^2\frac{S-as}{\rho^2},\qquad
 D_-X_{10}=\I Q^2\frac{S+as}{(\rho^*)^2}.
\label{eq:offdiagderiv}
\end{equation}
Direct multiplication of \(A\) and \(X\) yields
\begin{align}
 (AX)_{00}&=Q\frac{M\rho^*-Q^2}{\rho^2},&
 (AX)_{01}&=-\I Q^2\frac{S-as}{\rho^2},\nonumber\\
 (AX)_{10}&=\I Q^2\frac{S+as}{(\rho^*)^2},&
 (AX)_{11}&=Q\frac{M\rho-Q^2}{(\rho^*)^2}.
\label{eq:AXentries}
\end{align}
Thus, entry by entry,
\begin{equation}
 D_-X=AX.
\label{eq:DX}
\end{equation}

\begin{lemma}[Differentiating the inverse]
\label{lem:inverse}
On \(\Uext\),
\begin{equation}
 D_-(X^{-1})=-X^{-1}A.
\label{eq:DinvX}
\end{equation}
Moreover,
\begin{equation}
 D_+(X^\dagger)=X^\dagger A,\qquad
 D_+\bigl((X^\dagger)^{-1}\bigr)=-A(X^\dagger)^{-1}.
\label{eq:Dstar}
\end{equation}
\end{lemma}

\begin{proof}
Differentiate \(X^{-1}X=I\) with \(D_-\).  The product rule and
\eqref{eq:DX} give
\[
 D_-(X^{-1})X+X^{-1}AX=0.
\]
Right multiplication by \(X^{-1}\) proves \eqref{eq:DinvX}.  Complex
conjugation changes \(\I\) to \(-\I\); hence conjugating the first
identity \eqref{eq:DX} and taking the transpose gives
\(D_+(X^\dagger)=X^\dagger A^\dagger\).
Use \(A^\dagger=A\) from \eqref{eq:trace-hermitian}.  Differentiate
\((X^\dagger)^{-1}X^\dagger=I\) with \(D_+\), and cancel as above, to
obtain the second identity in \eqref{eq:Dstar}.
\end{proof}

\begin{proposition}[Residual transport through the integrating factor]
\label{prop:transport}
For every smooth two-component field \(v\), one has
\begin{align}
 D_-(Xv)-A(Xv)&=X(D_-v),\label{eq:transport-minus}\\
 D_+\bigl((X^\dagger)^{-1}v\bigr)
 +A\bigl((X^\dagger)^{-1}v\bigr)
 &=(X^\dagger)^{-1}(D_+v).\label{eq:transport-plus}
\end{align}
\end{proposition}

\begin{proof}
The product rule gives
\[
 D_-(Xv)=(D_-X)v+X(D_-v)=AXv+X(D_-v),
\]
which is \eqref{eq:transport-minus}.  Set \(Y=(X^\dagger)^{-1}\).
Then \eqref{eq:Dstar} gives
\[
 D_+(Yv)=(D_+Y)v+Y(D_+v)=-AYv+Y(D_+v),
\]
which proves \eqref{eq:transport-plus}.
\end{proof}

\begin{corollary}[Homogeneous solution correspondence]
\label{cor:homogeneous}
At every \(x\in\Uext\),
\[
 D_-v=0\quad\Longleftrightarrow\quad D_-(Xv)-A(Xv)=0,
\]
and
\[
 D_+v=0\quad\Longleftrightarrow\quad
 D_+\bigl((X^\dagger)^{-1}v\bigr)
 +A\bigl((X^\dagger)^{-1}v\bigr)=0.
\]
\end{corollary}

\begin{proof}
The forward directions follow from Proposition~\ref{prop:transport}.
For the reverse directions, \(X\) and \(X^\dagger\) are invertible by
\eqref{eq:detX}.
\end{proof}

\section{The transformed system and its equivalence}
\label{sec:transformation}

\subsection{Definition of system (7)}

For each spin sign, set
\begin{equation}
 z_+=\adj(X)y_+,\qquad z_-=X^\dagger y_-,
\label{eq:z}
\end{equation}
and define
\begin{equation}
 C^{(\varsigma)}
 =(\det X)X^{-1}
 \left(c^{(\varsigma)}I-B^{(\varsigma)}\right)
 (X^\dagger)^{-1}.
\label{eq:C}
\end{equation}
System (7) is
\begin{equation}
 D_-z_+=C^{(\varsigma)}z_-,\qquad
 D_+z_-=\adj(C^{(\varsigma)})z_+.
\label{eq:S7}
\end{equation}
All coefficients are smooth on \(\Uext\).  The constant determinant
identity makes both transformations in \eqref{eq:z} invertible.

\subsection{Adjugate identity for the coupling matrix}

For any \(2\times2\) matrix \(T\),
\(\adj(T)=\tr(T)I-T\).  Since \(\tr B^{(\varsigma)}=0\),
\begin{equation}
 \adj(c^{(\varsigma)}I-B^{(\varsigma)})
 =c^{(\varsigma)}I+B^{(\varsigma)}.
\label{eq:adjcB}
\end{equation}
Since \(X^{-1}=(\det X)^{-1}\adj(X)\) and
\((X^\dagger)^{-1}=(\det X)^{-1}\adj(X^\dagger)\), the definition
\eqref{eq:C} can be written
\[
 C^{(\varsigma)}
 =(\det X)^{-1}\adj(X)
 (c^{(\varsigma)}I-B^{(\varsigma)})\adj(X^\dagger).
\]
For \(2\times2\) matrices, \(\adj(UV)=\adj(V)\adj(U)\),
\(\adj(\adj(U))=U\), and \(\adj(\lambda U)=\lambda\adj(U)\).
Consequently,
\begin{align}
 \adj(C^{(\varsigma)})
 &=(\det X)^{-1}X^\dagger
 \adj(c^{(\varsigma)}I-B^{(\varsigma)})X\nonumber\\
 &=(\det X)^{-1}X^\dagger
 \left(c^{(\varsigma)}I+B^{(\varsigma)}\right)X.
\label{eq:adjC}
\end{align}
The scalar is real because \(\det X=Q^2-M^2\).

\begin{proposition}[Forward implication: system (4) implies system (7)]
\label{prop:forward}
If smooth \(y_+,y_-\) satisfy \eqref{eq:S4minus}--\eqref{eq:S4plus}
at \(x\in\Uext\), then \(z_+,z_-\) defined by \eqref{eq:z} satisfy
\eqref{eq:S7} at \(x\).
\end{proposition}

\begin{proof}
From \(\adj(X)=(\det X)X^{-1}\), \eqref{eq:DinvX}, and the constancy
of \(\det X\),
\[
 D_-\adj(X)=-(\adj X)A.
\]
The product rule therefore gives
\[
 D_-(\adj(X)y_+)
 =-\adj(X)Ay_++\adj(X)D_-y_+.
\]
Use \eqref{eq:S4minus}, namely
\(D_-y_+=Ay_++(cI-B)y_-\).  The connection terms cancel:
\[
 D_-z_+=\adj(X)(cI-B)y_-.
\]
Since \(y_-=(X^\dagger)^{-1}z_-\), and \(\adj(X)=(\det X)X^{-1}\),
this is the first equation of \eqref{eq:S7}.

For the other equation, use \(D_+(X^\dagger)=X^\dagger A\):
\[
 D_+(X^\dagger y_-)
 =X^\dagger Ay_-+X^\dagger D_+y_-.
\]
Equation \eqref{eq:S4plus} gives
\(D_+y_-=-Ay_-+(cI+B)y_+\), so
\[
 D_+z_-=X^\dagger(cI+B)y_+.
\]
Finally, \eqref{eq:adjC} and \(z_+=\adj(X)y_+\) show that
\[
 \adj(C)z_+
 =(\det X)^{-1}X^\dagger(cI+B)X\adj(X)y_+
 =X^\dagger(cI+B)y_+.
\]
Here \(X\adj(X)=(\det X)I\).  This proves the second equation.
\end{proof}

\begin{proposition}[Reverse implication: system (7) implies system (4)]
\label{prop:reverse}
Suppose smooth \(z_+,z_-\) satisfy \eqref{eq:S7} at \(x\in\Uext\).  Define
\begin{equation}
 y_+=(\det X)^{-1}Xz_+,\qquad
 y_-=(X^\dagger)^{-1}z_-.
\label{eq:inverse-z}
\end{equation}
Then \(y_+,y_-\) satisfy \eqref{eq:S4minus}--\eqref{eq:S4plus} at \(x\).
\end{proposition}

\begin{proof}
Apply \eqref{eq:transport-minus} with \(v=z_+\).  Since \(\det X\) is
constant,
\[
 D_-y_+-Ay_+
 =(\det X)^{-1}X(D_-z_+)
 =(\det X)^{-1}XCz_-.
\]
Substitute the definition of \(C\):
\[
 (\det X)^{-1}XC
 =(\det X)^{-1}X(\det X)X^{-1}(cI-B)(X^\dagger)^{-1}
 =(cI-B)(X^\dagger)^{-1}.
\]
As \(y_-=(X^\dagger)^{-1}z_-\), the first equation of system (4)
follows.

Next apply \eqref{eq:transport-plus} with \(v=z_-\):
\[
 D_+y_-+Ay_-=(X^\dagger)^{-1}(D_+z_-)
 =(X^\dagger)^{-1}\adj(C)z_+.
\]
Use \eqref{eq:adjC} and \(y_+=(\det X)^{-1}Xz_+\):
\[
 (X^\dagger)^{-1}\adj(C)z_+
 =(\det X)^{-1}(cI+B)Xz_+
 =(cI+B)y_+.
\]
This is the second equation of system (4).  All cancellations are
permitted by \(\det X\neq0\).
\end{proof}

\begin{corollary}[Two-sided equivalence]
\label{cor:47}
For either spin sign and every \(x\in\Uext\), system (4) holds at \(x\)
if and only if system (7) holds at \(x\) for the transformed fields.
\end{corollary}

\begin{proof}
Proposition~\ref{prop:forward} proves one direction and
Proposition~\ref{prop:reverse} the other.  The transformations round-trip:
\[
 (\det X)^{-1}X\adj(X)y_+=y_+,\qquad
 (X^\dagger)^{-1}X^\dagger y_-=y_-.
\]
\end{proof}

\section{Return to the four-field variables}
\label{sec:completion}

Proposition~\ref{prop:residual} identifies the vanishing of the original
four residuals with system (4), separately for positive and negative spin.
Corollary~\ref{cor:47} identifies system (4) with system (7).  Their
composition proves Theorem~\ref{thm:main}.

The chiral map itself has a pointwise algebraic inverse.  For positive
spin,
\begin{align}
 \psi_2&=\frac{y_+^0+y_-^0}{2P_+},&
 \psi_1&=\frac{\sqrt\Delta\,(y_+^0-y_-^0)}{2\I P_+},\label{eq:inverse-positive-a}\\
 k&=-\frac{y_+^1+y_-^1}{4QP_+},&
 s&=-\frac{\sqrt\Delta\,(y_+^1-y_-^1)}{4\I QP_+}.
\label{eq:inverse-positive-b}
\end{align}
For negative spin,
\begin{align}
 \psi_{-2}&=\frac{y_+^0+y_-^0}{2P_-},&
 \psi_{-1}&=-\frac{y_+^0-y_-^0}{2\I\sqrt\Delta\,P_-},\label{eq:inverse-negative-a}\\
 n&=\frac{y_+^1+y_-^1}{4QP_-},&
 \ell&=-\frac{y_+^1-y_-^1}{4\I Q\sqrt\Delta\,P_-}.
\label{eq:inverse-negative-b}
\end{align}
Substitution verifies both round trips.  Nonvanishing of \(Q\), \(P_\pm\),
and \(\sqrt\Delta\) is sufficient for every denominator in these formulas.
The regularity class is preserved in both directions because all
coefficients and inverse coefficients are smooth on \(\Uext\).

\subsection{Why the angular derivative is essential}

The angular part of \(D_-\) is essential for this particular matrix.
If
\[
 D_{\mathrm{bad}}=\sqrt\Delta\,\partial_r,
\]
then for \(a\neq0\) one has
\begin{equation}
 \bigl(D_{\mathrm{bad}}X_{ij}\bigr)_{i,j}\neq AX
\label{eq:bad-operator}
\end{equation}
at every point of \(\Uext\).  Indeed,
\[
 \bigl(D_{\mathrm{bad}}X-AX\bigr)_{01}
 =\I\partial_\theta X_{01}
 =-\frac{\I aQ^2\sin\theta}{\rho^2}\neq0.
\]
All factors in the numerator are nonzero under \(a\neq0\) and the
standing hypotheses.  As a concrete example, \((M,a,Q)=(2,1,1)\) is
subextremal and \((r,\theta)=(4,\pi/2)\) lies in its exterior.

\section{Scope and relation to the perturbation problem}
\label{sec:scope}

The result proved here is local and algebraic in the following precise
sense: at each exterior point, the two finite systems of differential
equations have equivalent residuals under the displayed smooth,
invertible transformations.  The quantification over arbitrary smooth
fields shows that the equivalence is an identity between differential
systems, not an assertion that a particular solution exists.

Several mathematically substantial questions are outside this result.
We do not derive the original equations from the linearized
Einstein--Maxwell equations; analyze horizon or axis regularity; impose
ingoing or outgoing boundary conditions; construct angular eigenfunctions;
derive a radial spectral problem; or prove mode stability.  Nor do we
claim that \(C^{(\varsigma)}\) separates as a sum of a radial matrix and
an angular matrix.  Explicit candidate coefficients for such a
decomposition can be written down, but proving that decomposition is not
part of the local equivalence established here.

In the physical mode ansatz, the azimuthal number is integral.  The
identities proved here are algebraic in that parameter and therefore hold
for every real \(m\), which is a slightly broader statement than the
physical mode convention.  Likewise, \(\omega\) is real throughout this
local result; no claim about complex-frequency modes follows.

\section{Proof provenance and reproducibility}
\label{sec:provenance}

An independently checkable proof archive \cite{Wang2026} accompanies this mathematical
exposition and records the definitions and derivations of the identities
used above.  The archived source is version 0.1.0 at
\href{https://doi.org/10.5281/zenodo.23684720}{doi:10.5281/zenodo.23684720}.
The archive establishes the local equivalence stated in
Theorem~\ref{thm:main}; it does not establish the analytic claims excluded
in Section~\ref{sec:scope}.

\appendix
\section{A component calculation for the chiral residuals}
\label{app:component}

\subsection{One positive-spin component}

The residual identities in Proposition~\ref{prop:residual} hold for
arbitrary fields, before any field equation is imposed.  A worked
component makes explicit how both derivative and zero-order terms are
checked.  Consider positive spin and set
\[
 S=\sqrt\Delta,\qquad
 \lambda=\frac{5\Delta'}{4S}-\frac{\I}{2}\cot\theta,\qquad
 H=\frac{3\I a\sin\theta}{\rho^*},\qquad
 T=\frac3{\rho^*}.
\]
Write \(g_\varepsilon=y_\varepsilon/P_+\), so that
\[
 g_\varepsilon^0=\psi_2+\varepsilon\I\psi_1/S,\qquad
 g_\varepsilon^1=-2Q(k+\varepsilon\I s/S).
\]
Since \(D_-P_+=P_+\lambda\),
\begin{align}
 \frac{\cE_{-,0}^{(+)}}{P_+}
 ={}&D_-g_+^0+\lambda g_+^0
       -A_{01}g_+^1-(c-\I b)g_-^0+B_{01}g_-^1\nonumber\\
 ={}&(S\partial_r-\I\partial_\theta)\psi_2
       +(\I\partial_r+S^{-1}\partial_\theta)\psi_1
       +\alpha\psi_2+\beta\psi_1\nonumber\\
 &\quad+2Q(A_{01}-B_{01})k
       +\frac{2\I Q}{S}(A_{01}+B_{01})s,
\label{eq:component-expand}
\end{align}
where
\begin{equation}
 \alpha=\lambda-c+\I b,\qquad
 \beta=\frac{\I}{S}(\lambda+c-\I b)
                 -\frac{\I\Delta'}{2\Delta}.
\label{eq:alphabeta}
\end{equation}
Here \(b,c,B\) have spin sign \(+1\).  Substituting their definitions,
and using \(S^2=\Delta\), gives the scalar coefficient checks
\begin{align}
 \alpha
 &=-\I(\widehat Q+2\cot\theta-H)
     +\frac{-\I K+2\Delta'-\Delta T}{S},
 \label{eq:alpha-check}\\
 \beta
 &=-\frac K\Delta+\I T
     +\frac{-\widehat Q-\cot\theta+H}{S},
 \label{eq:beta-check}\\
 2Q(A_{01}-B_{01})&=-2\I F_+,\qquad
 \frac{2\I Q}{S}(A_{01}+B_{01})=-\frac{2F_-}{S}.
 \label{eq:coupling-check}
\end{align}
The expression in \eqref{eq:component-expand}, with these coefficients,
is exactly \(-\I\Res_{1a}+\Res_{1b}/S\), upon expanding
\eqref{eq:R1a} and \eqref{eq:R1b}.  This proves the first component of
\eqref{eq:positive-residuals} without assuming that either residual
vanishes.

\subsection{A uniform calculation for both spins and both chiralities}
\label{app:uniform}

The following calculation verifies the other components with no implicit
change-of-sign convention.  Fix
\(\varsigma,\varepsilon\in\{+1,-1\}\), let \(S=\sqrt\Delta\), and use
the following four fields:
\[
\begin{array}{c|cccc}
 \varsigma & u&v&w&t\\ \hline
 +1&\psi_2&\psi_1&k&s\\
 -1&\psi_{-2}&-\psi_{-1}&-n&\ell
\end{array}
\]
Define
\[
 h=S^{-\varsigma},\qquad
 P=\Delta^{\varsigma+1/4}\sqrt{\sin\theta},\qquad
 g_\varepsilon=\mat{u+\varepsilon\I h v\\
                         -2Q(w+\varepsilon\I h t)}.
\]
Thus \(y_\varepsilon=Pg_\varepsilon\) is exactly the chiral field in
the relevant sector.  In the rest of this calculation \(b,c,B\) have
spin sign \(\varsigma\); \(H=3\I a\sin\theta/\rho^*\),
\(T=3/\rho^*\), and \(Z=\rho/(\rho^*)^2\).
Introduce the scalar differential operators
\begin{align}
 \mathsf L&=\partial_\theta+\varsigma\widehat Q+2\cot\theta-H,
 &\mathsf J&=\partial_\theta-\varsigma\widehat Q-\cot\theta+H,
 \label{eq:uniform-angular}\\
 \mathsf U&=S^{1-\varsigma}
 \left(\partial_r+\frac{\varsigma\I K}{\Delta}
                 +\frac{\varsigma-1}{2}\frac{\Delta'}{\Delta}+T\right),
 \label{eq:uniform-U}\\
 \mathsf V&=S^{1+\varsigma}
 \left(\partial_r-\frac{\varsigma\I K}{\Delta}
                 +(1+\varsigma)\frac{\Delta'}{\Delta}-T\right).
 \label{eq:uniform-V}
\end{align}
For \(\varsigma=+1\), the last two are
\(\D_0+T\) and \(\Delta(\D_2^\dagger-T)\); for
\(\varsigma=-1\), they are
\(\Delta(\D_{-1}^\dagger+T)\) and \(\D_0-T\).
Consequently the original four residuals in either sector are
\begin{align}
 \Res_a&=\mathsf L u-\mathsf U v+2F_+w,
 &\Res_b&=\mathsf V u+\mathsf J v-2F_-t,\nonumber\\
 \Res_c&=\mathsf U t-\mathsf J w-Zu,
 &\Res_d&=\mathsf V w+\mathsf L t-2Zv.
 \label{eq:uniform-original}
\end{align}
They are \(\Res_{1a},\ldots,\Res_{1d}\) for positive spin and
\(\Res^-_{1a},\ldots,\Res^-_{1d}\) for negative spin.

Put
\[
 \mathfrak d_\varepsilon=S\partial_r-\varepsilon\I\partial_\theta,
 \qquad
 \lambda_\varepsilon
   =\frac{\mathfrak d_\varepsilon P}{P}
   =\left(\varsigma+\frac14\right)\frac{\Delta'}{S}
         -\frac{\varepsilon\I}{2}\cot\theta.
\]
In particular,
\(\mathfrak d_\varepsilon h=-\varsigma h\Delta'/(2S)\).
For every smooth scalar field \(f\), expansion of the operators gives
the following four identities:
\begin{align}
 (\mathfrak d_\varepsilon+\lambda_\varepsilon-c
                    +\varepsilon\I b)f
   &=(-\varepsilon\I\mathsf L+h\mathsf V)f,
 \label{eq:uniform-op0}\\
 \mathfrak d_\varepsilon(\varepsilon\I h f)
   +\varepsilon\I h(\lambda_\varepsilon+c-\varepsilon\I b)f
   &=(\varepsilon\I\mathsf U+h\mathsf J)f,
 \label{eq:uniform-op1}\\
 (\mathfrak d_\varepsilon+\lambda_\varepsilon-c
                    -\varepsilon\I b)f
   &=(-\varepsilon\I\mathsf J+h\mathsf V)f,
 \label{eq:uniform-op2}\\
 \mathfrak d_\varepsilon(\varepsilon\I h f)
   +\varepsilon\I h(\lambda_\varepsilon+c+\varepsilon\I b)f
   &=(\varepsilon\I\mathsf U+h\mathsf L)f.
 \label{eq:uniform-op3}
\end{align}
To check these without suppressing the zero-order terms, observe that
\[
 hS^{1+\varsigma}=S,\qquad hS=S^{1-\varsigma}.
\]
The zeroth-order coefficient on either side of
\eqref{eq:uniform-op0} is
\[
 -\varepsilon\I(\varsigma\widehat Q+2\cot\theta-H)
 +\frac{-\varsigma\I K+(1+\varsigma)\Delta'-\Delta T}{S},
\]
and that on either side of \eqref{eq:uniform-op1} is
\[
 h\left[
 -\frac{\varepsilon\varsigma K}{S}
 +\frac{\varepsilon\I(\varsigma-1)\Delta'}{2S}
 +\varepsilon\I ST-\varsigma\widehat Q-\cot\theta+H
 \right].
\]
For \eqref{eq:uniform-op2}, replace the angular coefficient in
\eqref{eq:uniform-op0} by
\(-\varepsilon\I(-\varsigma\widehat Q-\cot\theta+H)\).
For \eqref{eq:uniform-op3}, replace the final angular expression in
\eqref{eq:uniform-op1} by
\(\varsigma\widehat Q+2\cot\theta-H\).
These four comparisons verify every derivative and multiplication term.
The off-diagonal coefficients require only
\begin{align}
 2Q(A_{01}-B_{01})&=-2\I F_+,
 &2\I Q(A_{01}+B_{01})&=-2F_-,\nonumber\\
 \frac{A_{10}-B_{10}}{2Q}&=-\I Z,
 &\frac{\I(A_{10}+B_{10})}{2Q}&=-2Z.
 \label{eq:uniform-couplings}
\end{align}
Each follows immediately from \eqref{eq:A}, \eqref{eq:B}, and
\eqref{eq:Fpm}.

Now let \(\mathfrak E_\varepsilon\) be the chiral residual belonging to
\(y_\varepsilon\):
\[
 \mathfrak E_\varepsilon
 =\mathfrak d_\varepsilon y_\varepsilon
   -\varepsilon Ay_\varepsilon
   -(cI-\varepsilon B)y_{-\varepsilon}.
\]
It is \(\cE_-\) when \(\varepsilon=+1\) and \(\cE_+\) when
\(\varepsilon=-1\).  Apply
\eqref{eq:uniform-op0}--\eqref{eq:uniform-op3} to \(u,v,w,t\),
respectively, and use \eqref{eq:uniform-couplings}.  The first component
becomes
\begin{align*}
 P^{-1}\mathfrak E_\varepsilon^0
 &=(-\varepsilon\I\mathsf L+h\mathsf V)u
   +(\varepsilon\I\mathsf U+h\mathsf J)v
   -2\varepsilon\I F_+w-2hF_-t\\
 &=-\varepsilon\I\Res_a+h\Res_b.
\end{align*}
Dividing the second component by \(-2QP\) gives
\begin{align*}
 (-2QP)^{-1}\mathfrak E_\varepsilon^1
 &=(-\varepsilon\I\mathsf J+h\mathsf V)w
   +(\varepsilon\I\mathsf U+h\mathsf L)t
   -\varepsilon\I Zu-2hZv\\
 &=\varepsilon\I\Res_c+h\Res_d.
\end{align*}
We have therefore proved the single vector identity
\begin{equation}
 P^{-1}\mathfrak E_\varepsilon
 =\mat{-\varepsilon\I\Res_a+h\Res_b\\
       -2Q(\varepsilon\I\Res_c+h\Res_d)}.
 \label{eq:uniform-final}
\end{equation}
Taking the two values of \(\varsigma\) and of \(\varepsilon\)
recovers all eight component identities in
\eqref{eq:positive-residuals}--\eqref{eq:negative-residuals-plus}.

\begin{thebibliography}{9}
\bibitem{Hintz2026}
P.~Hintz,
\emph{Perturbations of Kerr--Newman black holes: separation and mode
stability},
arXiv:2609.33661v1 [gr-qc] (2026).
\newblock \url{https://arxiv.org/abs/2609.33661}.

\bibitem{Wang2026}
H.~Wang,
\emph{Hintz2026: A Lean 4 Formalization of the Local Equivalence of
Kerr--Newman Perturbation Systems (1) and (7)},
Zenodo, version 0.1.0 (2026).
\newblock DOI: \url{https://doi.org/10.5281/zenodo.23684720}.
\end{thebibliography}

\end{document}
